\lhead[ \leftmark ]{Informatik, Informatik \& Gesellschaft}
\rhead[Informatik]{Cyril Wendl, \today, Winterthur}

\title{Bewertungsraster Factsheet\\ \large Informatik \& Gesellschaft}

\maketitle

\section*{Hinweise zum Factsheet}

\textbf{Bewertungskriterien}:
\begin{itemize}
    \item (2 Pt) Inhaltlich korrekte und umfassende Zusammenfassung des Kapitels
    \item (1 Pt) Mindestens eine selbst erstellte \textbf{Illustration} (Bild, Diagramm, Grafik, usw.), die Ihr Factsheet sinnvoll ergänzt.
    \item (1 Pt) Formale Kriterien
    \begin{itemize}
        \item Abgabeformat: PDF
        \item Dateiname: ``\texttt{[Klassenname] - K[Kapitelnummer] [Kapitelname] Factsheet.pdf}'' (z.B. ``\texttt{2d - K2 Daten Factsheet.pdf}'')
        \item 3-4 A4-Seiten
        \item Korrekte Sprache, vollständige Sätze
        \item Schriftgrösse: 11 Punkte
        \item Vor- und Nachnamen aller Gruppenmitglieder im PDF enthalten
    \end{itemize}
    \item (1 Pt) Formulieren Sie am Ende Ihrer Zusammenfassung zwei interessante, reichhaltige Diskussions-Fragen (die nicht einfach mit ja / nein beantwortet werden können) oder provokativ-pointierte Aussagen, die zur Diskussion anregen. Versuchen Sie, Fragen zu formulieren, die Bezug zu Ihrer eigenen Lebenswelt (Interessen, Hobbies etc.) haben.
\end{itemize}


Das Fact Sheet enthält zudem spannende Fragestellungen, welche die Grundlage für die Gruppendiskussionen bilden

Geben Sie im Anhang (zählt nicht zu den 3-4 Seiten) an, wer für welche(n) Unterabschnitt(e) zuständig war.

\vspace{1em}

Das Factsheet ist spätestens eine Woche vor dem Vortrag abzugeben.


\newpage
\part*{2e}
\newpage

% --- Bewertung Factsheet: Justiz und Algorithmen (2e-K3JustizFactsheet.pdf) ---

\section*{Bewertung: Justiz und Algorithmen (2e)}

\textbf{Autorinnen: Ladina, Maha, Minna, Salome, Aysima}

\subsection*{Factsheet}

\begin{itemize}
    \item \textbf{Inhaltliche Zusammenfassung (1.5/2 Punkte):} Die wichtigsten Aspekte werden behandelt, die Beispiele sind gut gewählt. Die Rechercheaufträge (z.B. Funktionsweise von Random Forests, Beispiele ausserhalb der Justiz, aktuelle Datenschutzregeln) werden ebenfalls gut und genau bearbeitet. Die Begriffe Sensitivität/Spezifität werden leider nicht erwähnt.
    \item \textbf{Illustrationen (.75/1 Punkt):} Es sind zwei passende Visualisierungen enthalten. Diese sind jeoch leider kaum lesbar, da die Auflösung zu niedrig und das Bild zu klein sind. Die Problematik der falsch-positiven und falsch-negativen Ergebnisse wird angesprochen, wäre aber durch eine eigene Visualisierung (z.B. Kreise) noch anschaulicher darstellbar.
    \item \textbf{Formale Kriterien (1/1 Punkt):} Alle formalen Vorgaben sind erfüllt.
    \item \textbf{Diskussionsfragen (1/1 Punkt):} Gute, sinnvolle und offene Fragen.
\end{itemize}

\textbf{Gesamtpunktzahl: 4.25/5 Punkten}

\subsection*{Vortrag}
\footnotesize
\begin{tabular}{p{10cm}p{6cm}}
\midrule\midrule
\textbf{Inhalt: Gesamteindruck} & \\
Das Thema scheint nicht verstanden worden zu sein. (0 Punkte) & \\
Das Thema scheint nur teilweise verstanden worden zu sein. (2 Punkte) &  \\
Das Thema wurde gut und im Detail verstanden. (\underline{4 Punkte}) & \emoji{check-mark-button} Sie haben eine vollständige, spannende und anschauliche Präsentation gehalten. Bravo!\\
\midrule\midrule
\textbf{Inhalt: Qualität des Vortrags} & \\
Inhalte beschrieben, aber nicht logisch erklärt. (0 Punkte) & \\
Erklärung teilweise unklar. (2 Punkte) & \emoji{check-mark-button} ``Einzelner Entscheidungsbaum nicht sehr gut'' $\rightarrow$ eventuell genauer erklären, weshalb (\textit{overfitting})\\
Alle Inhalte klar und logisch erklärt. (4 Punkte) &  \\
\midrule\midrule
\textbf{Inhalt: Fachsprache und Fachbegriffe} & \\
Mehrmals falsche Fachsprache. (0 Punkte) & \\
Mindestens ein Fehler bei der Fachsprache. (1 Punkt) & \\
Fachsprache immer korrekt. (2 Punkte) & \emoji{check-mark-button} \\
\midrule\midrule
\textbf{Inhalt: Einbettung des Themas} & \\
Nicht oder falsch eingebettet. (0 Punkte) & \\
Vollständig richtig eingebettet. (2 Punkte) & \emoji{check-mark-button} \\
\midrule\midrule
\textbf{Inhalt: Aufbau und Struktur} & \\
Struktur führte zu Verwirrung. (0 Punkte) & \\
Struktur förderlich für das Verständnis. (2 Punkte) & \emoji{check-mark-button} \\
\midrule\midrule
\textbf{Inhalt: Vollständigkeit} & \\
Ein zentraler Inhalt fehlt. (-4 Punkte) & \\
Alle zentralen Inhalte vollständig. (0 Punkte) & \emoji{check-mark-button} \\
\midrule\midrule
\textbf{Form / Medien: Eigenständigkeit} & \\
Medien kopiert oder keine verwendet. (-4 Punkte) & \\
Eigenständige Medien. (2 Punkte) & \\
Neue Perspektive durch Medien. (4 Punkte) & \emoji{check-mark-button} Theater-Einlage top gemacht, sehr unterhaltsam und kreativ! Bravo!\\
\midrule\midrule
\textbf{Form / Medien: Übersichtlichkeit} & \\
Medien unübersichtlich oder nicht aussagekräftig. (0 Punkte) & \\
Übersichtliche und aussagekräftige Medien. (1 Punkt) & \\
Herausragende Medien. (2 Punkte) & \emoji{check-mark-button} Präsentation sorgfältig und anschaulich gestaltet.\\
\midrule\midrule
\textbf{Form / Medien: Anwendungen} & \\
Keine Beispiele / Anwendungen. (0 Punkte) & \\
Beispiele wie im Unterricht. (2 Punkte) & \\
Kreative, eigene Beispiele. (4 Punkte) & \emoji{check-mark-button} s. oben \\
\midrule\midrule
\textbf{Form / Medien: Interaktion mit der Klasse} & \\
Keine Einbindung oder gestellt. (0 Punkte) & \\
Sinnvolle Einbindung und Anregung zum Mitdenken. (2 Punkte) & \emoji{check-mark-button} Gute, komplexe und anspruchsvolle Fragen \\
\midrule\midrule
\textbf{Auftreten und Vortragsweise} & \\
Abgelesen, monoton oder unpassend. (0 Punkte) & \\
Frei gesprochen, positives Auftreten. (2 Punkte) & \emoji{check-mark-button} Angenehmes auftreten.\\
\midrule\midrule
\textbf{Selbstorganisation} & \\
Termin/Thema vergessen, keine Nachfrage. (-4 Punkte) & \\
Datum notiert, vorgängig informiert. (0 Punkte) & \emoji{check-mark-button} \\
\midrule\midrule
\textbf{Vortragsdauer} & \\
Weniger als 30 oder weniger als 20 Minuten. (-2 Punkte) & \\
Zwischen 20 und 30 Minuten. (0 Punkte) & \emoji{check-mark-button}  \\
\midrule\midrule
\normalsize
\end{tabular}
\textbf{Gesamtpunktzahl:} \underline{24 / 26 Punkte}

\textbf{Note}
\begin{itemize}
    \item[\textbf{Factsheet}] 5.25
    \item[\textbf{Vortrag}] 5.6
    \item[\textbf{Gesamtnote}]  \underline{5.4}
\end{itemize}


\newpage

% --- Bewertung Factsheet: Medizin und KI (2e-K4-Medizin-Factsheet.pdf) ---

\section*{Bewertung Factsheet: Medizin und KI}

\textbf{Autor:innen: Miguel Berger, Andrii Marchenko, Anna Greminger, Anna Graf}

\begin{itemize}
    \item \textbf{Inhaltliche Zusammenfassung (1.5/2 Punkte):} Die Funktionsweise neuronaler Netze und die Anwendung in der Medizin werden verständlich erklärt. Die Rechercheaufträge (z.B. CNN ausserhalb der Medizin, Beispiele aus Schweizer Spitälern, Datenschutzprobleme) werden angesprochen, aber nicht immer mit konkreten, aktuellen Beispielen oder kritischer Reflexion vertieft.
    \item \textbf{Illustrationen (1/1 Punkt):} Eigene Visualisierungen sind vorhanden und sinnvoll.
    \item \textbf{Formale Kriterien (1/1 Punkt):} Alle Vorgaben sind erfüllt.
    \item \textbf{Diskussionsfragen (0.5/1 Punkt):} Die Fragen sind relevant, könnten aber noch provokativer oder lebensweltlicher formuliert sein.
\end{itemize}

\textbf{Gesamtpunktzahl: 4/5 Punkten}

\newpage

% --- Bewertung Factsheet: Autonomes Fahren und Bayes (2e-K5-Autos-Factsheet.pdf) ---

\section*{Bewertung Factsheet: Autonomes Fahren und Bayes}

\textbf{Autor:innen: Eli Hess, Flynn Sallefranque, Lenny Schnetzer, Sharon Onesimus}

\begin{itemize}
    \item \textbf{Inhaltliche Zusammenfassung (1.5/2 Punkte):} Die Zusammenfassung ist verständlich und deckt die wichtigsten technischen und ethischen Aspekte ab. Die Bayes-Erklärung ist anschaulich, aber die Verbindung zu den konkreten Rechercheaufträgen (z.B. weitere Bayes-Anwendungen, aktuelle Forschung zu autonomen Autos) bleibt teilweise oberflächlich. Einige Abschnitte wirken zu erzählerisch und weniger analytisch.
    \item \textbf{Illustrationen (0.5/1 Punkt):} Es gibt anschauliche Beispiele, aber die Visualisierungen sind eher beschreibend als analytisch. Ein selbst erstelltes, datenbasiertes Diagramm (z.B. Unfallstatistik, Entscheidungsbaum) fehlt.
    \item \textbf{Formale Kriterien (1/1 Punkt):} Sprache, Länge und Format stimmen, alle Namen sind enthalten.
    \item \textbf{Diskussionsfragen (1/1 Punkt):} Die Fragen sind offen und regen zur Diskussion an.
\end{itemize}

\textbf{Gesamtpunktzahl: 4/5 Punkten}

\newpage

% --- Bewertung Factsheet: Kriminalität, Algorithmen und Gesichtserkennung (2e-K6Kriminalitaet-Factsheet.pdf) ---

\section*{Bewertung Factsheet: Kriminalität, Algorithmen und Gesichtserkennung}

\textbf{Autor:innen: [Namen bitte ergänzen]}

\begin{itemize}
    \item \textbf{Inhaltliche Zusammenfassung (1.5/2 Punkte):} Die Entwicklung von Algorithmen in der Kriminalitätsbekämpfung wird gut dargestellt. Die Rechercheaufträge (z.B. Polizei-IT-Systeme in der Schweiz, Social Credit Score in China) werden nur teilweise oder oberflächlich behandelt. Die kritische Reflexion zu Diskriminierung und Datenschutz ist vorhanden, aber könnte noch vertieft werden.
    \item \textbf{Illustrationen (1/1 Punkt):} Visualisierungen sind vorhanden und passend.
    \item \textbf{Formale Kriterien (1/1 Punkt):} Alle Vorgaben sind erfüllt.
    \item \textbf{Diskussionsfragen (0.5/1 Punkt):} Die Fragen sind offen, könnten aber noch stärker auf aktuelle gesellschaftliche Debatten eingehen.
\end{itemize}

\textbf{Gesamtpunktzahl: 4/5 Punkten}

\newpage

